\chapter{Stapedectomy}

\section*{What Is This Surgery?}
Stapedectomy is a specialized microsurgical procedure performed by otologic surgeons to address conductive hearing loss primarily caused by otosclerosis, a condition characterized by abnormal bone growth around the stapes bone in the middle ear. The surgery involves the removal of the fixed stapes bone, either partially (stapedotomy) or completely (stapedectomy), from the oval window, which is the membrane-covered opening to the inner ear. The primary goal of this surgery is to restore the efficient transmission of sound vibrations from the tympanic membrane through the ossicular chain to the perilymphatic fluid of the cochlea. By replacing the non-functional stapes with a prosthetic device, typically a piston, the surgery aims to re-establish the mechanical connection between the incus and the inner ear, significantly improving the patient's hearing acuity and overall quality of life.

\section*{Alternative Names}
\begin{itemize}
  \item Stapedotomy
  \item Stapes surgery
  \item Stapes prosthesis placement
  \item Ossicular chain reconstruction (specifically involving the stapes)
  \item Stapes fenestration
  \item Oval window fenestration
\end{itemize}

\section*{Relevant Surgical Anatomy}
The middle ear is an air-filled cavity located within the temporal bone, housing three tiny bones known as the ossicles: the malleus, incus, and stapes. These ossicles form a chain that transmits sound vibrations from the tympanic membrane (eardrum) to the oval window of the inner ear. The stapes, the smallest bone in the human body, consists of a head, two crura, and a footplate. The footplate fits into the oval window, sealed by an annular ligament, and is responsible for transmitting vibrations to the perilymph of the cochlea. 

Key anatomical structures in relation to the stapes include:

\begin{itemize}
  \item **Facial Nerve (Cranial Nerve VII):** This nerve runs through the middle ear in the Fallopian canal, often dehiscent in parts, and is located superior to the oval window. Injury to this nerve can lead to facial paralysis.
  \item **Chorda Tympani Nerve:** A branch of the facial nerve that carries taste sensations, it crosses the middle ear cavity medial to the malleus and incus, making it vulnerable to injury during surgery.
  \item **Promontory:** A bulge formed by the basal turn of the cochlea, located inferior to the oval window.
  \item **Arterial Supply:** The middle ear receives blood supply from branches of the external carotid artery, including the anterior tympanic artery and stylomastoid artery.
  \item **Venous Drainage:** Venous drainage parallels the arterial supply, primarily through the pterygoid plexus.
  \item **Lymphatic Drainage:** Lymphatic vessels follow the venous pathways to regional lymph nodes.
\end{itemize}

Understanding these anatomical relationships is crucial for minimizing complications during surgery and ensuring optimal outcomes.

\section*{Surgical Principle \& Rationale}
The fundamental principle of stapedectomy is to address the pathology of otosclerosis, characterized by abnormal bone remodeling in the otic capsule, leading to the fixation of the stapes footplate within the oval window. This fixation prevents the stapes from vibrating freely in response to sound waves, impeding the transmission of mechanical energy to the perilymphatic fluid of the inner ear and resulting in conductive hearing loss. The rationale for surgical intervention is to bypass this mechanical block and restore the impedance matching mechanism of the middle ear.

During the procedure, the surgeon aims to create a new, mobile pathway for sound conduction. By carefully excising the fixed stapes superstructure and creating a small opening (fenestra) in the footplate (stapedotomy) or removing the entire footplate (stapedectomy), direct access to the inner ear fluids is established. A precisely sized prosthetic piston, typically made of biocompatible materials such as Teflon or titanium, is then inserted. One end of the prosthesis is crimped onto the long process of the incus, while the other end extends through the fenestra into the inner ear, often sealed with a small tissue graft (e.g., fat or fascia). This re-establishes the vibratory connection from the tympanic membrane and ossicular chain directly to the inner ear fluids, allowing sound to be perceived effectively and improving hearing acuity.

\section*{History \& Surgical Evolution}
The history of stapedectomy dates back to the late 19th and early 20th centuries, with early attempts to address stapes fixation through procedures like fenestration surgery. Julius Lempert pioneered this technique in the 1930s, creating a new opening in the lateral semicircular canal to bypass the ossicular chain. However, these early procedures had limited success and were associated with significant complications, including persistent vertigo and sensorineural hearing loss.

The modern era of stapes surgery began in 1956 when Dr. John J. Shea Jr. performed the first successful total stapedectomy, removing the entire stapes and replacing it with a Teflon prosthesis and vein graft. This groundbreaking technique revolutionized the treatment of otosclerosis, providing a direct and effective solution for conductive hearing loss.

In the following decades, the procedure evolved significantly. Surgeons like Dr. Howard P. House and Dr. James L. Sheehy refined the technique, contributing to its widespread adoption and standardization. A major advancement was the introduction of stapedotomy by Dr. John Perkins and Dr. Ugo Fisch in the 1970s. This modification involves creating only a small fenestra in the stapes footplate, which reduces the risk of perilymph fistula and sensorineural hearing loss by minimizing trauma to the inner ear. The introduction of lasers (e.g., CO2, KTP) and microdrills for precise fenestration further enhanced safety and precision, solidifying stapedectomy and stapedotomy as the gold standard for treating otosclerosis.

\section*{Clinical Indications}
Stapedectomy is indicated for patients presenting with clinically significant conductive hearing loss primarily due to otosclerosis. Specific indications include:

\begin{itemize}
  \item Significant conductive hearing loss due to clinically confirmed otosclerosis, typically identified through audiometric evaluation and clinical examination.
  \item An air-bone gap of 20 dB or greater on pure tone audiometry, indicating a mechanical impediment to sound conduction that is amenable to surgical correction.
  \item Good cochlear reserve, meaning the sensorineural component of hearing loss is not severe, suggesting a high potential for significant hearing improvement post-surgery.
  \item Patient preference for surgical correction over long-term hearing aid use, after thorough counseling on risks and benefits.
  \item Progressive hearing loss attributed to otosclerosis, impacting daily communication, social interaction, and overall quality of life.
  \item Presence of bothersome tinnitus associated with otosclerosis that may improve or resolve with restoration of hearing.
  \item Unilateral or bilateral otosclerosis, with surgery typically performed on the poorer-hearing ear first to assess outcomes before considering the contralateral ear.
  \item Absence of active middle ear infection, external ear infection, or other contraindications that would preclude safe surgical intervention.
\end{itemize}

\section*{Clinical Positioning}
Stapedectomy is considered the primary surgical treatment for conductive hearing loss caused by otosclerosis. For eligible patients, it is often the first-line intervention offered after a thorough diagnostic workup confirms the diagnosis and assesses the extent of hearing loss. While hearing aids can effectively manage the symptoms of otosclerosis by amplifying sound, stapedectomy offers a potential anatomical and physiological correction of the underlying pathology, often leading to a significant and lasting improvement in natural hearing.

In some cases, it may be considered a secondary procedure, such as a revision stapedectomy for a failed or displaced prosthesis from a previous surgery, or if a patient initially opted for hearing aids but later desires a more definitive surgical solution. It is not typically used as a salvage procedure for other forms of hearing loss or middle ear pathology unless otosclerosis is the underlying cause. Its role is highly specific to the diagnosis of otosclerosis, distinguishing it from other middle ear surgeries like tympanoplasty or ossiculoplasty for different etiologies.

\section*{Surgical Technique \& Procedure}
Stapedectomy is a meticulous microsurgical procedure typically performed under general anesthesia, though local anesthesia with sedation can also be used for cooperative patients. The patient is positioned supine with the head turned to the side, allowing optimal access to the ear canal and middle ear structures.

The procedure involves the following key steps:

\begin{enumerate}
  \item \textbf{Access and Exposure:} An ear speculum is inserted into the external auditory canal. A tympanomeatal flap, which is a small incision in the ear canal skin and tympanic membrane, is carefully elevated using micro-instruments to expose the middle ear space and the ossicular chain (malleus, incus, stapes).
  
  \item \textbf{Diagnosis Confirmation \& Ossicular Disarticulation:} The surgeon confirms the diagnosis of stapes fixation by gently palpating the ossicles. The incudostapedial joint, connecting the incus to the stapes, is then carefully disarticulated. The stapedial tendon, if present, is cut.
  
  \item \textbf{Stapes Superstructure Removal:} The stapes superstructure (the crura and head) is carefully fractured and removed, leaving only the fixed stapes footplate in the oval window.
  
  \item \textbf{Fenestration of the Footplate:} A small opening, or fenestra, is created in the fixed stapes footplate within the oval window. This critical step can be performed using a microdrill, a laser (e.g., CO2 or KTP laser), or a fine pick. The size of the fenestra is crucial, typically 0.6-0.8 mm in diameter, to minimize inner ear trauma.
  
  \item \textbf{Graft Placement \& Prosthesis Insertion:} A small tissue graft, often a piece of fat from the earlobe or fascia from the temporalis muscle, is then placed over the fenestra to seal the inner ear and prevent perilymph leakage. A prosthetic piston, chosen for its appropriate length and diameter, is then inserted. One end of the piston is passed through the tissue graft into the fenestra, while the other end is crimped securely onto the long process of the incus.
  
  \item \textbf{Prosthesis Verification \& Closure:} The mobility of the prosthesis and the entire ossicular chain is then gently checked to ensure proper sound transmission. The tympanomeatal flap is carefully repositioned back into place. Absorbable packing material, such as Gelfoam or cotton, may be placed in the ear canal to hold the flap in position and absorb any minor bleeding. No external skin incision or sutures are typically required for primary stapedectomy.
\end{enumerate}

\section*{Preoperative Workup \& Preparation}
Thorough preoperative evaluation is crucial for a successful stapedectomy, ensuring patient safety and optimal outcomes. This typically includes:

\begin{itemize}
  \item \textbf{Audiometric Evaluation:} Comprehensive pure tone audiometry, speech audiometry, tympanometry, and acoustic reflex testing are performed to confirm conductive hearing loss, quantify the air-bone gap, assess cochlear function, and rule out other pathologies.
  
  \item \textbf{Imaging Studies:} A high-resolution computed tomography (CT) scan of the temporal bones may be ordered to confirm the diagnosis of otosclerosis (e.g., lucency around the oval window), rule out other middle ear pathologies (e.g., cholesteatoma, ossicular discontinuity), and identify any anatomical variations (e.g., facial nerve dehiscence, persistent stapedial artery) that could complicate surgery.
  
  \item \textbf{Medical Clearance:} A complete medical history and physical examination are conducted. Patients may require cardiac clearance or other specialist consultations (e.g., endocrinologist for diabetes) to ensure they are medically fit for anesthesia and surgery, especially if they have significant comorbidities.
  
  \item \textbf{Medication Review:} Patients are instructed to discontinue blood-thinning medications (e.g., aspirin, NSAIDs, warfarin, novel oral anticoagulants) for a specified period (typically 7-10 days) before surgery to minimize bleeding risk. Herbal supplements with anticoagulant properties should also be stopped.
  
  \item \textbf{NPO Status:} Patients must be nil per os (NPO) for a minimum of 6-8 hours prior to surgery to prevent aspiration during general anesthesia.
  
  \item \textbf{Antibiotics:} A single dose of prophylactic broad-spectrum antibiotics (e.g., cefazolin) is often administered intravenously just before the incision to reduce the risk of postoperative infection.
  
  \item \textbf{Informed Consent:} A detailed discussion with the surgeon regarding the procedure, potential benefits, risks, alternatives (e.g., hearing aids), expected outcomes, and recovery process. The patient must fully understand and provide informed consent.
\end{itemize}

\section*{Contraindications \& Precautions}
\subsection*{Absolute Contraindications}
\begin{itemize}
  \item \textbf{Only Hearing Ear:} Surgery on the only functional hearing ear is generally contraindicated due to the catastrophic risk of total sensorineural hearing loss, unless it is the only viable option for bilateral severe otosclerosis and the patient fully understands and accepts this profound risk.
  
  \item \textbf{Active Middle Ear or External Ear Infection:} Any active infection (acute otitis media, otitis externa, cholesteatoma) must be thoroughly treated and completely resolved before surgery to prevent serious complications like labyrinthitis, meningitis, or prosthesis infection.
  
  \item \textbf{Retrocochlear Pathology:} Evidence of a retrocochlear lesion (e.g., acoustic neuroma, vestibular schwannoma) that could mimic otosclerosis symptoms or complicate the surgical field.
  
  \item \textbf{Poor General Health:} Patients with severe systemic diseases, uncontrolled chronic conditions (e.g., severe heart failure, uncontrolled hypertension), or those who are medically unstable for anesthesia and surgery.
\end{itemize}

\subsection*{Relative Contraindications \& Precautions}
\begin{itemize}
  \item \textbf{Perforated Tympanic Membrane:} A pre-existing tympanic membrane perforation should ideally be repaired prior to stapedectomy, as it significantly increases the risk of infection, perilymph fistula, and poor healing.
  
  \item \textbf{Non-aerated Middle Ear:} Conditions like chronic otitis media with effusion, Eustachian tube dysfunction, or significant middle ear scarring can compromise middle ear healing, prosthesis function, and long-term success.
  
  \item \textbf{Méni\`{e}re's Disease:} Patients with pre-existing Méni\`{e}re's disease may experience worsening vertigo or hearing loss post-operatively due to inner ear irritation.
  
  \item \textbf{Vertigo as a Primary Symptom:} While otosclerosis can cause dizziness, severe pre-operative vertigo may indicate a higher risk of persistent or worsened vestibular symptoms after surgery, requiring careful counseling.
  
  \item \textbf{Pregnancy:} Surgery is generally deferred until after pregnancy due to potential hormonal influences on otosclerosis progression and the inherent risks associated with anesthesia and medications during gestation.
  
  \item \textbf{Uncontrolled Diabetes:} Poorly controlled diabetes mellitus can impair wound healing, increase susceptibility to infection, and potentially affect microvascular supply to the inner ear.
  
  \item \textbf{Occupations with Pressure Changes:} Patients whose occupations or hobbies involve rapid or significant pressure changes (e.g., divers, pilots, frequent air travelers) should be thoroughly counseled on the increased risks of perilymph fistula, prosthesis displacement, or barotrauma.
\end{itemize}

\section*{Complications \& Risks}
While stapedectomy is generally safe and effective, like any surgical procedure, it carries potential risks and complications. These can be broadly categorized as general surgical risks and procedure-specific risks:

\begin{itemize}
  \item \textbf{General Surgical Risks:}
    \begin{itemize}
      \item \textbf{Bleeding:} Minor bleeding from the ear canal is common; significant hemorrhage requiring intervention is rare.
      \item \textbf{Infection:} Middle ear infection (otitis media) or external ear infection (otitis externa) can occur, potentially leading to more serious inner ear infections.
      \item \textbf{Anesthesia Complications:} Adverse reactions to anesthetic agents, respiratory depression, cardiovascular events, or aspiration pneumonia.
    \end{itemize}
  
  \item \textbf{Procedure-Specific Risks:}
    \begin{itemize}
      \item \textbf{Sensorineural Hearing Loss:} This is the most serious complication, ranging from partial high-frequency loss to profound, irreversible "dead ear" (0.5-2\% risk), often due to inner ear trauma or labyrinthitis.
      \item \textbf{Dizziness or Vertigo:} Transient dizziness and imbalance are common immediately post-op. Persistent or severe vertigo can occur due to inner ear irritation, perilymph fistula, or labyrinthine damage.
      \item \textbf{Tinnitus:} Pre-existing tinnitus may worsen, or new tinnitus may develop, though often it improves with successful hearing restoration.
      \item \textbf{Facial Nerve Injury:} The facial nerve runs close to the surgical field. Injury can cause temporary or permanent facial weakness or paralysis (rare, <1\%), ranging from mild paresis to complete paralysis.
      \item \textbf{Perilymph Fistula:} Leakage of inner ear fluid through the oval window, leading to fluctuating hearing loss, vertigo, and tinnitus. This may require prompt revision surgery to seal the leak.
      \item \textbf{Prosthesis Displacement or Extrusion:} The prosthetic piston may shift out of position, become dislodged, or extrude, leading to recurrent conductive hearing loss and requiring revision surgery.
      \item \textbf{Taste Disturbance:} Injury to the chorda tympani nerve, which crosses the middle ear, can cause altered taste sensation (dysgeusia) or loss of taste (ageusia) on one side of the tongue, often temporary but can be permanent.
      \item \textbf{Tympanic Membrane Perforation:} A hole in the eardrum can occur during flap elevation, manipulation, or repositioning, potentially requiring myringoplasty.
      \item \textbf{Ossicular Chain Discontinuity:} Damage to the incus or malleus during surgery, requiring further reconstruction or affecting prosthesis stability.
      \item \textbf{Failure to Improve Hearing:} Despite technically successful surgery, some patients may not achieve the desired hearing improvement due to underlying cochlear issues or inadequate sound transmission.
      \item \textbf{Revision Surgery:} Required in a small percentage of cases (5-10\%) due to complications (e.g., fistula, prosthesis displacement) or inadequate hearing improvement.
    \end{itemize}
\end{itemize}

\section*{Postoperative Recovery Timeline}
Immediately after stapedectomy, the patient is transferred to the post-anesthesia care unit (PACU) for close monitoring of vital signs, pain, and any signs of dizziness or nausea. Mild dizziness, some ear pain, and a sensation of ear fullness are common during this period. The ear canal is typically packed with absorbable material, which may cause muffled hearing. Most patients are discharged home the same day or the following morning, depending on their recovery from anesthesia and overall stability.

During the first 24-48 hours, patients are advised to rest, keep their head elevated, and take prescribed pain medication. They should strictly avoid sudden head movements, straining (e.g., during bowel movements), lifting heavy objects, and blowing their nose forcefully. Sneezing should be done with the mouth open to equalize pressure and prevent inner ear barotrauma. The operated ear must be kept absolutely dry, avoiding water entry into the ear canal during showering or bathing. Over the first 1-2 weeks, the ear canal packing may remain in place, and hearing may initially feel worse due to swelling and the packing. Patients are gradually allowed to resume light activities but should continue to avoid strenuous exercise, air travel, and activities that involve rapid pressure changes. Full hearing improvement is often gradual, becoming noticeable weeks to months after surgery as swelling subsides and the inner ear adapts to the new sound conduction pathway.

\section*{Surgical Findings \& Pathology}
During stapedectomy, the surgeon meticulously documents intraoperative findings. Key observations include the degree of stapes footplate fixation, which can range from mild stiffness to complete immobility. The presence of any anatomical variations, such as a dehiscent facial nerve canal overlying the oval window or a persistent stapedial artery, is noted as these can complicate the procedure. The condition of the incus and malleus, the integrity of the round window reflex, and the presence of any perilymphatic gush upon fenestration are also important findings.

If a portion of the stapes footplate or superstructure is removed, it is typically sent for histopathological examination. The pathology report will confirm the diagnosis of otosclerosis, characterized by abnormal bone remodeling within the otic capsule, often showing areas of active bone resorption and new bone formation (spongiosis) or dense, sclerotic bone (sclerosis). This microscopic confirmation validates the clinical diagnosis and provides definitive evidence of the underlying disease process responsible for the conductive hearing loss.

\section*{Expected Postoperative Course}
A normal, successful postoperative course following stapedectomy is characterized by a progressive improvement in hearing and resolution of initial surgical symptoms. Mild ear pain and discomfort are expected for a few days, managed with oral analgesics, and typically subside within the first week. Initial muffled hearing due to ear canal packing and middle ear swelling gradually improves as the packing is removed and swelling resolves, usually over 2-4 weeks. Dizziness or vertigo, if present immediately after surgery, should be transient and resolve within days to a few weeks.

Patients can expect to gradually resume normal daily activities, with full return to non-strenuous work within 1-2 weeks. The tympanomeatal flap heals without external sutures, and the ear canal returns to its normal appearance. The primary indicator of success is the significant improvement in hearing, with the air-bone gap closing to within 10 dB on audiometry, typically assessed at 4-6 weeks post-op. Any pre-existing tinnitus often improves or resolves as hearing is restored. The long-term course involves stable hearing improvement, allowing patients to communicate more effectively and often discontinue hearing aid use.

\section*{Postoperative Care \& Follow-up}
Postoperative care for stapedectomy is crucial for optimal healing and long-term success.

\begin{itemize}
  \item \textbf{Postoperative Visits:} The first follow-up visit is typically 1-2 weeks after surgery for removal of any ear canal packing, initial wound check, and assessment of early healing. A second visit, usually 4-6 weeks post-op, includes a formal postoperative audiogram to objectively assess hearing improvement and confirm air-bone gap closure.
  
  \item \textbf{Activity Restrictions:} Patients are advised to avoid strenuous physical activity, heavy lifting, nose blowing, and straining for at least 2-4 weeks. Air travel and activities involving rapid pressure changes (e.g., diving, skydiving) should be avoided indefinitely to prevent perilymph fistula or prosthesis displacement.
  
  \item \textbf{Ear Care:} The operated ear must be kept dry for several weeks. Patients are instructed on how to protect the ear during showering (e.g., using a cotton ball with petroleum jelly).
  
  \item \textbf{Medications:} Oral antibiotics may be prescribed for a short course post-op, along with pain relievers and anti-nausea medication as needed.
  
  \item \textbf{Vestibular Rehabilitation:} If persistent dizziness or imbalance occurs beyond the initial recovery period, referral to vestibular physical therapy may be recommended to help the inner ear adapt.
  
  \item \textbf{Warning Signs:} Patients are educated to immediately contact their surgeon if they experience sudden, severe hearing loss, persistent or worsening vertigo, new onset of facial weakness or paralysis, fever, purulent ear discharge, or severe ear pain.
\end{itemize}

Long-term follow-up may involve periodic audiograms to monitor hearing stability.

\section*{Epidemiology}
Otosclerosis, the primary indication for stapedectomy, affects approximately 0.3\% to 1\% of the general population, with clinical manifestations leading to hearing loss occurring in about 0.1\% of individuals. It is more prevalent in Caucasian populations and is significantly less common in individuals of African or Asian descent. The condition typically presents in young to middle-aged adults, most commonly between the ages of 20 and 40, though it can manifest earlier or later in life. There is a notable female predominance, with women being affected about twice as often as men. The disease often progresses or worsens during pregnancy, suggesting a hormonal influence. While otosclerosis is the most frequent cause of conductive hearing loss in adults without a history of chronic otitis media, the overall mortality associated with stapedectomy is extremely rare, primarily related to general anesthesia risks. Morbidity, however, can include the risks outlined previously, with sensorineural hearing loss being the most significant potential complication, occurring in a small percentage of cases.

\section*{Outcomes \& Prognosis}
The outcomes of stapedectomy for otosclerosis are generally excellent, with high success rates in restoring hearing. Approximately 90-95\% of patients achieve closure of the air-bone gap to within 10 dB on postoperative audiometry, leading to significant subjective improvement in hearing and overall quality of life. Functional outcomes are typically very good, with most patients experiencing improved communication abilities, reduced listening effort, and often a reduced reliance on or complete discontinuation of hearing amplification. The prognosis for long-term hearing stability after a successful stapedectomy is favorable, with many patients maintaining their improved hearing for decades. True recurrence of otosclerosis at the operated oval window is rare, but new foci of otosclerosis can develop in other parts of the otic capsule or in the contralateral ear, potentially requiring future intervention. Prognostic factors influencing success include the surgeon's experience, the extent and type of otosclerosis, the absence of other inner ear pathologies, and the patient's overall health. While revision surgery may be necessary in a small percentage of cases (5-10\%) due to complications like prosthesis displacement or perilymph fistula, the overall long-term outlook for hearing improvement and patient satisfaction is overwhelmingly positive.

\section*{References}
\begin{enumerate}
  \item MedlinePlus. Stapedectomy. U.S. National Library of Medicine; 2024. Available from: \url{https://medlineplus.gov/ency/article/002998.htm}
  
  \item Flint PW, Haughey BH, Lund VJ, et al., editors. Cummings Otolaryngology: Head and Neck Surgery. 7th ed. Elsevier; 2021.
  
  \item American Academy of Otolaryngology--Head and Neck Surgery Foundation. Clinical Practice Guideline: Otosclerosis. 2023.
  
  \item Shea JJ Jr. Fenestration of the oval window. Ann Otol Rhinol Laryngol. 1956;65(4):1021-1033.
\end{enumerate}

\clearpage